\IfFileExists{stacks-project.cls}{%
\documentclass{stacks-project}
}{%
\documentclass{amsart}
}

% For dealing with references we use the comment environment
\usepackage{verbatim}
\newenvironment{reference}{\comment}{\endcomment}
%\newenvironment{reference}{}{}
\newenvironment{slogan}{\comment}{\endcomment}
\newenvironment{history}{\comment}{\endcomment}

% For commutative diagrams we use Xy-pic
\usepackage[all]{xy}

% We use 2cell for 2-commutative diagrams.
\xyoption{2cell}
\UseAllTwocells

% We use multicol for the list of chapters between chapters
\usepackage{multicol}

% This is generally recommended for better output
\usepackage{lmodern}
\usepackage[T1]{fontenc}

% For cross-file-references
\usepackage{xr-hyper}

% Package for hypertext links:
\usepackage{hyperref}

% For any local file, say "hello.tex" you want to link to please
% use \externaldocument[hello-]{hello}
\externaldocument[introduction-]{introduction}
\externaldocument[conventions-]{conventions}
\externaldocument[sets-]{sets}
\externaldocument[categories-]{categories}
\externaldocument[topology-]{topology}
\externaldocument[sheaves-]{sheaves}
\externaldocument[sites-]{sites}
\externaldocument[stacks-]{stacks}
\externaldocument[fields-]{fields}
\externaldocument[algebra-]{algebra}
\externaldocument[brauer-]{brauer}
\externaldocument[homology-]{homology}
\externaldocument[derived-]{derived}
\externaldocument[simplicial-]{simplicial}
\externaldocument[more-algebra-]{more-algebra}
\externaldocument[smoothing-]{smoothing}
\externaldocument[modules-]{modules}
\externaldocument[sites-modules-]{sites-modules}
\externaldocument[injectives-]{injectives}
\externaldocument[cohomology-]{cohomology}
\externaldocument[sites-cohomology-]{sites-cohomology}
\externaldocument[dga-]{dga}
\externaldocument[dpa-]{dpa}
\externaldocument[sdga-]{sdga}
\externaldocument[hypercovering-]{hypercovering}
\externaldocument[schemes-]{schemes}
\externaldocument[constructions-]{constructions}
\externaldocument[properties-]{properties}
\externaldocument[morphisms-]{morphisms}
\externaldocument[coherent-]{coherent}
\externaldocument[divisors-]{divisors}
\externaldocument[limits-]{limits}
\externaldocument[varieties-]{varieties}
\externaldocument[topologies-]{topologies}
\externaldocument[descent-]{descent}
\externaldocument[perfect-]{perfect}
\externaldocument[more-morphisms-]{more-morphisms}
\externaldocument[flat-]{flat}
\externaldocument[groupoids-]{groupoids}
\externaldocument[more-groupoids-]{more-groupoids}
\externaldocument[etale-]{etale}
\externaldocument[chow-]{chow}
\externaldocument[intersection-]{intersection}
\externaldocument[pic-]{pic}
\externaldocument[weil-]{weil}
\externaldocument[adequate-]{adequate}
\externaldocument[dualizing-]{dualizing}
\externaldocument[duality-]{duality}
\externaldocument[discriminant-]{discriminant}
\externaldocument[derham-]{derham}
\externaldocument[local-cohomology-]{local-cohomology}
\externaldocument[algebraization-]{algebraization}
\externaldocument[curves-]{curves}
\externaldocument[resolve-]{resolve}
\externaldocument[models-]{models}
\externaldocument[functors-]{functors}
\externaldocument[equiv-]{equiv}
\externaldocument[pione-]{pione}
\externaldocument[etale-cohomology-]{etale-cohomology}
\externaldocument[proetale-]{proetale}
\externaldocument[relative-cycles-]{relative-cycles}
\externaldocument[more-etale-]{more-etale}
\externaldocument[trace-]{trace}
\externaldocument[crystalline-]{crystalline}
\externaldocument[spaces-]{spaces}
\externaldocument[spaces-properties-]{spaces-properties}
\externaldocument[spaces-morphisms-]{spaces-morphisms}
\externaldocument[decent-spaces-]{decent-spaces}
\externaldocument[spaces-cohomology-]{spaces-cohomology}
\externaldocument[spaces-limits-]{spaces-limits}
\externaldocument[spaces-divisors-]{spaces-divisors}
\externaldocument[spaces-over-fields-]{spaces-over-fields}
\externaldocument[spaces-topologies-]{spaces-topologies}
\externaldocument[spaces-descent-]{spaces-descent}
\externaldocument[spaces-perfect-]{spaces-perfect}
\externaldocument[spaces-more-morphisms-]{spaces-more-morphisms}
\externaldocument[spaces-flat-]{spaces-flat}
\externaldocument[spaces-groupoids-]{spaces-groupoids}
\externaldocument[spaces-more-groupoids-]{spaces-more-groupoids}
\externaldocument[bootstrap-]{bootstrap}
\externaldocument[spaces-pushouts-]{spaces-pushouts}
\externaldocument[spaces-chow-]{spaces-chow}
\externaldocument[groupoids-quotients-]{groupoids-quotients}
\externaldocument[spaces-more-cohomology-]{spaces-more-cohomology}
\externaldocument[spaces-simplicial-]{spaces-simplicial}
\externaldocument[spaces-duality-]{spaces-duality}
\externaldocument[formal-spaces-]{formal-spaces}
\externaldocument[restricted-]{restricted}
\externaldocument[spaces-resolve-]{spaces-resolve}
\externaldocument[formal-defos-]{formal-defos}
\externaldocument[defos-]{defos}
\externaldocument[cotangent-]{cotangent}
\externaldocument[examples-defos-]{examples-defos}
\externaldocument[algebraic-]{algebraic}
\externaldocument[examples-stacks-]{examples-stacks}
\externaldocument[stacks-sheaves-]{stacks-sheaves}
\externaldocument[criteria-]{criteria}
\externaldocument[artin-]{artin}
\externaldocument[quot-]{quot}
\externaldocument[stacks-properties-]{stacks-properties}
\externaldocument[stacks-morphisms-]{stacks-morphisms}
\externaldocument[stacks-limits-]{stacks-limits}
\externaldocument[stacks-cohomology-]{stacks-cohomology}
\externaldocument[stacks-perfect-]{stacks-perfect}
\externaldocument[stacks-introduction-]{stacks-introduction}
\externaldocument[stacks-more-morphisms-]{stacks-more-morphisms}
\externaldocument[stacks-geometry-]{stacks-geometry}
\externaldocument[moduli-]{moduli}
\externaldocument[moduli-curves-]{moduli-curves}
\externaldocument[examples-]{examples}
\externaldocument[exercises-]{exercises}
\externaldocument[guide-]{guide}
\externaldocument[desirables-]{desirables}
\externaldocument[coding-]{coding}
\externaldocument[obsolete-]{obsolete}
\externaldocument[fdl-]{fdl}
\externaldocument[index-]{index}

% Theorem environments.
%
\theoremstyle{plain}
\newtheorem{theorem}[subsection]{Theorem}
\newtheorem{proposition}[subsection]{Proposition}
\newtheorem{lemma}[subsection]{Lemma}

\theoremstyle{definition}
\newtheorem{definition}[subsection]{Definition}
\newtheorem{example}[subsection]{Example}
\newtheorem{exercise}[subsection]{Exercise}
\newtheorem{situation}[subsection]{Situation}

\theoremstyle{remark}
\newtheorem{remark}[subsection]{Remark}
\newtheorem{remarks}[subsection]{Remarks}

\numberwithin{equation}{subsection}

% Macros
%
\def\lim{\mathop{\mathrm{lim}}\nolimits}
\def\colim{\mathop{\mathrm{colim}}\nolimits}
\def\Spec{\mathop{\mathrm{Spec}}}
\def\Hom{\mathop{\mathrm{Hom}}\nolimits}
\def\Ext{\mathop{\mathrm{Ext}}\nolimits}
\def\SheafHom{\mathop{\mathcal{H}\!\mathit{om}}\nolimits}
\def\SheafExt{\mathop{\mathcal{E}\!\mathit{xt}}\nolimits}
\def\Sch{\mathit{Sch}}
\def\Mor{\mathop{\mathrm{Mor}}\nolimits}
\def\Ob{\mathop{\mathrm{Ob}}\nolimits}
\def\Sh{\mathop{\mathit{Sh}}\nolimits}
\def\NL{\mathop{N\!L}\nolimits}
\def\CH{\mathop{\mathrm{CH}}\nolimits}
\def\proetale{{pro\text{-}\acute{e}tale}}
\def\etale{{\acute{e}tale}}
\def\QCoh{\mathit{QCoh}}
\def\Ker{\mathop{\mathrm{Ker}}}
\def\Im{\mathop{\mathrm{Im}}}
\def\Coker{\mathop{\mathrm{Coker}}}
\def\Coim{\mathop{\mathrm{Coim}}}

% Boxtimes
%
\DeclareMathSymbol{\boxtimes}{\mathbin}{AMSa}{"02}

%
% Macros for moduli stacks/spaces
%
\def\QCohstack{\mathcal{QC}\!\mathit{oh}}
\def\Cohstack{\mathcal{C}\!\mathit{oh}}
\def\Spacesstack{\mathcal{S}\!\mathit{paces}}
\def\Quotfunctor{\mathrm{Quot}}
\def\Hilbfunctor{\mathrm{Hilb}}
\def\Curvesstack{\mathcal{C}\!\mathit{urves}}
\def\Polarizedstack{\mathcal{P}\!\mathit{olarized}}
\def\Complexesstack{\mathcal{C}\!\mathit{omplexes}}
% \Pic is the operator that assigns to X its picard group, usage \Pic(X)
% \Picardstack_{X/B} denotes the Picard stack of X over B
% \Picardfunctor_{X/B} denotes the Picard functor of X over B
\def\Pic{\mathop{\mathrm{Pic}}\nolimits}
\def\Picardstack{\mathcal{P}\!\mathit{ic}}
\def\Picardfunctor{\mathrm{Pic}}
\def\Deformationcategory{\mathcal{D}\!\mathit{ef}}
\usepackage{cleveref}
\crefname{originallemma}{original lemma}{original lemmas}
\Crefname{originallemma}{Original lemma}{Original lemmas}
\makeatletter
\newlabel{algebra-lemma-completion-Noetherian-Noetherian}{{0316}{source}{}{sourcepointer-1}{}}
\newlabel{algebra-lemma-completion-Noetherian-Noetherian@cref}{{[originallemma][0][]0316}{[1][0][]source}{}{}{}}
\makeatother
\crefname{originallemma}{original lemma}{original lemmas}
\Crefname{originallemma}{Original lemma}{Original lemmas}
\makeatletter
\newlabel{algebra-lemma-qc-open}{{00F6}{source}{}{sourcepointer-2}{}}
\newlabel{algebra-lemma-qc-open@cref}{{[originallemma][0][]00F6}{[1][0][]source}{}{}{}}
\makeatother
\crefname{originallemma}{original lemma}{original lemmas}
\Crefname{originallemma}{Original lemma}{Original lemmas}
\makeatletter
\newlabel{lemma-completion-affine-formal-is-affine-formal}{{0GVQ}{source}{}{sourcepointer-3}{}}
\newlabel{lemma-completion-affine-formal-is-affine-formal@cref}{{[originallemma][0][]0GVQ}{[1][0][]source}{}{}{}}
\makeatother
\crefname{originallemma}{original lemma}{original lemmas}
\Crefname{originallemma}{Original lemma}{Original lemmas}
\makeatletter
\newlabel{lemma-countably-indexed}{{0AN1}{source}{}{sourcepointer-4}{}}
\newlabel{lemma-countably-indexed@cref}{{[originallemma][0][]0AN1}{[1][0][]source}{}{}{}}
\makeatother
\crefname{originallemma}{original lemma}{original lemmas}
\Crefname{originallemma}{Original lemma}{Original lemmas}
\makeatletter
\newlabel{lemma-mcquillan-affine-formal-algebraic-space}{{0AIC}{source}{}{sourcepointer-5}{}}
\newlabel{lemma-mcquillan-affine-formal-algebraic-space@cref}{{[originallemma][0][]0AIC}{[1][0][]source}{}{}{}}
\makeatother
\crefname{originaldefinition}{original definition}{original definitions}
\Crefname{originaldefinition}{Original definition}{Original definitions}
\makeatletter
\newlabel{more-algebra-definition-topological-ring}{{07E8}{source}{}{sourcepointer-6}{}}
\newlabel{more-algebra-definition-topological-ring@cref}{{[originaldefinition][0][]07E8}{[1][0][]source}{}{}{}}
\makeatother
\crefname{originaldefinition}{original definition}{original definitions}
\Crefname{originaldefinition}{Original definition}{Original definitions}
\makeatletter
\newlabel{definition-weakly-adic}{{0GXE}{source}{}{sourcepointer-7}{}}
\newlabel{definition-weakly-adic@cref}{{[originaldefinition][0][]0GXE}{[1][0][]source}{}{}{}}
\makeatother
\title{Formal completion of an affine formal algebraic space along a closed subset keeps its type}
\author{The Stacks Project Authors}
\date{stacks-project:formal-spaces.tex:local-snapshot-20261012; GNU FDL 1.2 or later}
\begin{document}
\maketitle
\noindent\textit{Editorial note.} All delivered mathematical text is the author's own, copied from the pinned Stacks Project chapter file formal-spaces.tex; no statement or proof step was rewritten, reordered or removed.\par\medskip
\noindent\textit{Editorial note.} Part (4) of the proof ends with the author's 'Some details omitted.' It is reproduced verbatim and recorded as an obligation of the card; the parallel argument in part (3) is carried out in full and part (4) is explicitly said to be the same argument in the adic case.\par\medskip
\noindent\textit{Editorial note.} The proof displays three descriptions of the completion as a limit of quotients; all lie inside the counted span and are reproduced verbatim.\par\medskip
\noindent\textit{Editorial note.} The counted spans carry 6 navigation stubs; no bibliography citation occurs in the spans.\par\medskip
\noindent\textit{Editorial note.} The preamble span points at the chapter's own preamble.tex shipped with the frozen source, used verbatim as staged.\par\medskip
\noindent\textit{Allowed standard prerequisites (editorial).}\begin{itemize}
\item The formal completion of an affine formal algebraic space along a closed subset is affine. Same chapter, Lemma lemma-completion-affine-formal-is-affine-formal, invoked at proof lines 8122-8123. Gives part (1).
\item The second defining property of a McQuillan affine formal algebraic space. Same chapter, Lemma lemma-mcquillan-affine-formal-algebraic-space at proof lines 8124-8127 and Definition definition-types-affine-formal-algebraic-space at proof lines 8127-8128. Gives part (2).
\item Countably indexed affine formal algebraic spaces as formal spectra with a countable fundamental system of weak ideals of definition. Same chapter, Lemma lemma-countably-indexed, invoked at proof lines 8131-8134. Supplies the presentation used in the computation for part (3).
\item A closed subset with quasi-compact complement in an affine scheme is cut out by a finitely generated ideal. Algebra, Lemma algebra-lemma-qc-open, invoked at proof lines 8134-8137. Produces the finitely generated ideal whose lifts are used to describe the completion.
\item Completions of Noetherian rings. Algebra, Lemma algebra-lemma-completion-Noetherian-Noetherian, invoked at proof lines 8164-8165. Gives part (5) once the completion has been described as a limit of quotients.
\end{itemize}\medskip
\ifdefined\hypertarget\hypertarget{sourcepointer-1}{}\fi
\noindent\textit{Source reference.} Original lemma 0316: sources/stacks-chow/tags.txt, source lines 3973--3973. Pointer outside the counted span, invoked once for context; the load-bearing argument of the counted statement is the author's proof reproduced in full in the proof span.\par\medskip
\ifdefined\hypertarget\hypertarget{sourcepointer-2}{}\fi
\noindent\textit{Source reference.} Original lemma 00F6: sources/stacks-chow/tags.txt, source lines 560--560. Pointer outside the counted span, invoked once for context; the load-bearing argument of the counted statement is the author's proof reproduced in full in the proof span.\par\medskip
\ifdefined\hypertarget\hypertarget{sourcepointer-3}{}\fi
\noindent\textit{Source reference.} Original lemma 0GVQ: sources/stacks-formal-spaces/pinned-original-source.tex, source lines 7969--7969. Pointer outside the counted span, invoked once for context; the load-bearing argument of the counted statement is the author's proof reproduced in full in the proof span.\par\medskip
\ifdefined\hypertarget\hypertarget{sourcepointer-4}{}\fi
\noindent\textit{Source reference.} Original lemma 0AN1: sources/stacks-formal-spaces/pinned-original-source.tex, source lines 2493--2493. Pointer outside the counted span, invoked once for context; the load-bearing argument of the counted statement is the author's proof reproduced in full in the proof span.\par\medskip
\ifdefined\hypertarget\hypertarget{sourcepointer-5}{}\fi
\noindent\textit{Source reference.} Original lemma 0AIC: sources/stacks-formal-spaces/pinned-original-source.tex, source lines 2135--2135. Pointer outside the counted span, invoked once for context; the load-bearing argument of the counted statement is the author's proof reproduced in full in the proof span.\par\medskip
\ifdefined\hypertarget\hypertarget{sourcepointer-6}{}\fi
\noindent\textit{Source reference.} Original definition 07E8: sources/stacks-chow/tags.txt, source lines 9438--9438. Pointer outside the counted span, invoked once for context; the load-bearing argument of the counted statement is the author's proof reproduced in full in the proof span.\par\medskip
\ifdefined\hypertarget\hypertarget{sourcepointer-7}{}\fi
\noindent\textit{Source reference.} Original definition 0GXE: sources/stacks-formal-spaces/pinned-original-source.tex, source lines 1595--1595. Pointer outside the counted span, invoked once for context; the load-bearing argument of the counted statement is the author's proof reproduced in full in the proof span.\par\medskip
\begin{definition}
\label{definition-types-affine-formal-algebraic-space}
Let $S$ be a scheme. Let $X$ be an affine formal algebraic space over $S$.
We say $X$ is {\it McQuillan} if $X$ satisfies the equivalent conditions
of Lemma \ref{lemma-mcquillan-affine-formal-algebraic-space}. Let $A$
be the weakly admissible topological ring associated to $X$. We say
\begin{enumerate}
\item $X$ is {\it classical} if $X$ is McQuillan and $A$ is admissible
(More on Algebra, Definition \ref{more-algebra-definition-topological-ring}),
\item $X$ is {\it weakly adic} if $X$ is McQuillan and $A$ is weakly adic
(Definition \ref{definition-weakly-adic}),
\item $X$ is {\it adic} if $X$ is McQuillan and $A$ is adic
(More on Algebra, Definition \ref{more-algebra-definition-topological-ring}),
\item $X$ is {\it adic*} if $X$ is McQuillan, $A$ is adic, and $A$
has a finitely generated ideal of definition, and
\item $X$ is {\it Noetherian} if $X$ is McQuillan and $A$ is
both Noetherian and adic.
\end{enumerate}
\end{definition}

\begin{lemma}
\label{lemma-recompletion-affine-types}
Let $S$ be a scheme. Let $X$ be an affine formal algebraic space over $S$.
Let $T \subset X_{red}$ be a closed subset and let $X_{/T}$
be the formal completion of $X$ along $T$. Then
\begin{enumerate}
\item $X_{/T}$ is an affine formal algebraic space,
\item if $X$ is McQuillan, then $X_{/T}$ is McQuillan,
\item if $|X_{red}| \setminus T$ is quasi-compact and $X$
is countably indexed, then $X_{/T}$ is countably indexed,
\item if $|X_{red}| \setminus T$ is quasi-compact and $X$
is adic*, then $X_{/T}$ is adic*,
\item if $X$ is Noetherian, then $X_{/T}$ is Noetherian.
\end{enumerate}
\end{lemma}

\begin{proof}
Part (1) is Lemma \ref{lemma-completion-affine-formal-is-affine-formal}.
If $X$ is McQuillan, then $X = \text{Spf}(A)$ for some weakly admissible
topological ring $A$. Then $X_{/T} \to X \to \Spec(A)$ satisfies
property (2) of Lemma \ref{lemma-mcquillan-affine-formal-algebraic-space}
and hence $X_{/T}$ is McQuillan, see
Definition \ref{definition-types-affine-formal-algebraic-space}.

\medskip\noindent
Assume $X$ and $T$ are as in (3).
Then $X = \text{Spf}(A)$ where $A$ has a fundamental system
$A \supset I_1 \supset I_2 \supset I_3 \supset \ldots$
of weak ideals of definition, see Lemma \ref{lemma-countably-indexed}.
By Algebra, Lemma \ref{algebra-lemma-qc-open}
we can find a finitely generated ideal
$\overline{J} = (\overline{f}_1, \ldots, \overline{f}_r) \subset A/I_1$
such that $T$ is cut out by $\overline{J}$ inside $\Spec(A/I_1) = |X_{red}|$.
Choose $f_i \in A$ lifting $\overline{f}_i$.
If $Z = \Spec(B)$ is an affine scheme and $g : Z \to X$ is
a morphism with $g(Z) \subset T$ (set theoretically), then
$g^\sharp : A \to B$ factors through $A/I_n$ for some $n$
and $g^\sharp(f_i)$ is nilpotent in $B$ for each $i$. Thus
$J_{m, n} = (f_1, \ldots, f_r)^m + I_n$ maps to zero in $B$
for some $n, m \geq 1$. It follows that $X_{/T}$ is the formal spectrum of
$\lim_{n, m} A/J_{m, n}$ and hence countably indexed.
This proves (3).

\medskip\noindent
Proof of (4). Here the argument is the same as in (3).
However, here we may choose $I_n = I^n$ for some finitely generated ideal
$I \subset A$. Then it is clear that $X_{/T}$ is the formal spectrum of
$\lim A/J^n$ where $J = (f_1, \ldots, f_r) + I$. Some details omitted.

\medskip\noindent
Proof of (5). In this case $X_{red}$ is the spectrum of a Noetherian ring
and hence the assumption that $|X_{red}| \setminus T$ is quasi-compact
is satisfied. Thus as in the proof of (4) we see that
$X_{/T}$ is the spectrum of $\lim A/J^n$ which is a Noetherian
adic topological ring, see
Algebra, Lemma \ref{algebra-lemma-completion-Noetherian-Noetherian}.
\end{proof}

\end{document}
